\documentclass[11pt,a4paper]{article}
\usepackage[utf8]{inputenc}
\usepackage[T1]{fontenc}
\usepackage[brazil]{babel}
\usepackage[margin=2.2cm]{geometry}
\usepackage{amsmath}
\usepackage{parskip}

\newcommand{\slide}[2]{\subsection*{Slide #1 --- #2}}

\title{Roteiro de fala --- Segmentação 2D no Gray-Scott}
\author{Equipe Gray-Scott}
\date{Novembro de 2026}

\begin{document}
\maketitle

\noindent Tempo estimado: 6 a 8 minutos (cerca de 45 segundos por slide).

\slide{1}{Título}
Boa tarde. Vamos apresentar o andamento do nosso trabalho de PDI: a segmentação de padrões
no dataset Gray-Scott, do The Well.

\slide{2}{Objetivo e dataset}
A pergunta que guia o trabalho é: onde estão as estruturas e como elas evoluem no tempo?
Para responder, vamos segmentar os padrões de Turing do Gray-Scott, um sistema de
reação-difusão com duas concentrações químicas, $A$ e $B$. O experimento cruza dez
pré-processamentos, três redes (U-Net, SegFormer e Mask2Former) e dez pós-processamentos,
o que dá trezentas combinações. O dataset tem seis regimes, definidos pelos parâmetros $f$ e $k$,
e cada regime gera um padrão diferente: pontos, labirintos, espirais, entre outros.

\slide{3}{Situação do projeto}
Dividimos o trabalho em sete etapas. A auditoria do dataset está concluída. A definição do
ground truth está pronta para decisão, mas depende de uma resposta do professor. As etapas
de código, piloto e experimentos vêm na sequência. Como nossa máquina não tem GPU, os
treinos serão feitos no Colab ou no Kaggle.

\slide{4}{Auditoria: principais achados}
Antes de treinar, verificamos os dados. A divisão oficial tem 160 trajetórias de treino,
20 de validação e 20 de teste por regime. Lemos só as partes necessárias dos arquivos:
2,5 GB em vez de 154. Confirmamos as trajetórias estacionárias de Gliders e Spirals e
descobrimos que nelas o padrão desaparece por volta de 500 segundos, o que geraria máscaras
vazias. Também vimos que o instante inicial tem valores atípicos e que a documentação tem
um erro de sinal na equação de $B$.

\slide{5}{Tabela visual por instante}
Esta é a tabela visual pedida pelo professor: uma simulação de cada regime em seis instantes.
Ela mostra a evolução de cada padrão. Maze e Spots estabilizam cedo, enquanto Spirals e
Gliders continuam se movendo até o final. Essa diferença vai ser importante na definição do
ground truth.

\slide{6}{Ground truth: alternativas}
O ground truth é o gabarito: uma máscara que diz, pixel a pixel, onde há padrão. Como o
dataset não traz máscaras, nós as definimos a partir de $B$. Avaliamos três alternativas:
uma máscara binária por limiar, uma versão com três classes, que inclui a interface, e a
separação em componentes conexos. O limiar de 0,151 foi calculado pelo método de Otsu,
usando apenas os dados de treino.

\slide{7}{O ``frame médio'' representa o fluido?}
O professor sugeriu um ground truth baseado em um frame médio. Testamos isso medindo o
quanto uma máscara fixa coincide com os outros instantes. Nos regimes estáveis, como Maze
e Spots, funciona razoavelmente. Já em Gliders e Spirals o padrão se move e a coincidência
cai para menos de 20\%. Por isso, propomos que o frame médio defina o critério e que a
máscara seja gerada em cada instante.

\slide{8}{Decisões tomadas}
Já tomamos algumas decisões. A rede vai receber só $A$, porque a máscara vem de $B$: se a
rede visse $B$, ela apenas reaprenderia o limiar. Vamos excluir os frames sem padrão e o
instante inicial. Na condição ``sem pré'' usamos os dados brutos, e na ``sem pós'',
binarização com limiar fixo de 0,5.

\slide{9}{Perguntas e próximos passos}
Para seguir, precisamos de algumas respostas do professor, principalmente sobre o que é o
frame médio e se o ground truth deve ser um por regime. Com isso definido, implementamos o
código, rodamos um piloto, depois os 33 treinamentos, e por fim as tabelas e o artigo.
Obrigado.

\end{document}
